\chapter{Теория}
\label{ch:chap1}

  
\section{VLAN}

Virtual Local Area Network (VLAN) - технология, позволяющая на L2 коммутаторе создавать локальные сети. Делается
это путем добавления к пакету метки. Коммутатор не способен маршрутизировать пакеты между vlan. 
\section{Порты доступа}

Порт досутпа - порт, который предназначен для подключения оконечных устройств. Они добавляют метку к пакету при получении его
от оконечного устройства и снимают метку при отправке пакета оконечному устройству. Это необходимо потому, что оконечные
устройства не знают стандарт 802.1Q, который описывает метки. 

\section{Порты транка}
Транковый порт - порт, который способен передавать/принимать пакеты сразу с несколькими vlan. Такие порты настраиваются
между сетевыми устройствами.


Применение:

\begin{enumerate}
    \item Повышение производительности сети из-за снижения broadcast трафика
    \item Повышение безопасности. Устройства из одного vlan не могут получить доступ в другой vlan
    \item Удобство администрировани: устройства разделены на группы (по vlan), что делает тополгию сети более структурированной.
\end{enumerate}


\endinput